% Chapter 6: Online Learning and the Littlestone Dimension
% mistake_bounded = Profile B-light, regret_bounded = Profile C (Revelation),
% littlestone_characterization = Profile B (Load-Bearing).
% Centerpiece: the SOA algorithm and mistake tree diagram.
% Style: combinatorial and algorithmic, in contrast to Ch5's probabilistic flavor.

\chapter{Online Learning and the Littlestone Dimension}
\label{ch:online}

PAC learning can afford to be patient: nature draws training examples from a fixed distribution, the learner observes and generalizes, and the worst the world can do is an unlucky sample. That patience is a luxury. In online learning the luxury is revoked. There is no distribution to average over, no training phase to hide behind. At each round an adversary presents an instance, the learner commits to a label, and the adversary reveals the truth; every disagreement is a recorded mistake. The question is not ``what does the learner achieve in expectation?'' but ``how many mistakes can a clever adversary force, no matter what the learner does?''

This shift from distribution to adversary, from batch to sequential, from probabilistic to combinatorial, transforms the mathematics from the ground up. PAC learning is characterized by a number, the VC dimension, extracted through probabilistic concentration arguments. Online learning is characterized by a \emph{tree}, the mistake tree, whose depth the adversary can trace root-to-leaf to force exactly that many mistakes. The central result of this chapter, $\mathrm{Opt}(\mathcal{H}) = \Ldim(\mathcal{H})$, is proved by constructing both sides: an algorithm that limits mistakes to $\Ldim(\mathcal{H})$, and an adversary that forces at least $\Ldim(\mathcal{H})$. No tail probability appears anywhere; the entire argument is an explicit strategy meeting an explicit counter-strategy.

% ================================================================
% SECTION 1: THE ONLINE LEARNING GAME
% Profile G (Framework Bridge). Presented as a GAME, not a definition.
% ================================================================

\section{The Online Learning Game}
\label{sec:online-game}

The setup is a protocol between two players, a \emph{learner} and an \emph{adversary}, competing over a domain $X$, a label set $Y = \{0,1\}$, and a hypothesis class $\mathcal{H} \subseteq \{0,1\}^X$. One constraint ties the adversary's hands: \emph{realizability}, the requirement that some $h^* \in \mathcal{H}$ must be consistent with every label the adversary has assigned. Apart from that, the adversary is free to adapt, watching the learner's prior predictions before choosing the next instance $x_t$.

\begin{defn}[Online Learning Protocol]
\label{def:online-protocol}
The online learning game proceeds in rounds $t = 1, 2, 3, \ldots$\,:
\begin{enumerate}[leftmargin=*, itemsep=2pt]
\item The adversary selects an instance $x_t \in X$.
\item The learner observes $x_t$ and predicts a label $\hat{y}_t \in \{0,1\}$.
\item The adversary reveals the true label $y_t \in \{0,1\}$.
\item If $\hat{y}_t \neq y_t$, the learner incurs a \emph{mistake}.
\end{enumerate}
The game continues for as many rounds as the adversary chooses. The adversary must ensure that there exists $h^* \in \mathcal{H}$ with $h^*(x_t) = y_t$ for all $t$ (realizability).
\end{defn}

\begin{figure}[ht]
\centering
\begin{tikzpicture}[
    >={Stealth[length=2mm]},
    advbox/.style={draw, thick, rounded corners, fill=thmred!12, font=\small\bfseries,
      minimum width=2.5cm, minimum height=0.65cm},
    lrnbox/.style={draw, thick, rounded corners, fill=defblue!12, font=\small\bfseries,
      minimum width=2.5cm, minimum height=0.65cm},
    life/.style={semithick, densely dashed, gray!55},
    rstep/.style={-{Stealth[length=2mm]}, semithick},
    msg/.style={font=\scriptsize, fill=white, inner sep=1.5pt},
    val/.style={draw, semithick, rounded corners, align=center, font=\scriptsize,
      inner sep=3pt, minimum height=0.85cm, minimum width=1.9cm},
    vedge/.style={-{Stealth[length=2mm]}, sepgreen!55!black, semithick},
    sedge/.style={-{Stealth[length=2mm]}, edgeorange!80!black, semithick, densely dashed},
    hdl/.style={font=\tiny, text=black!70, align=center}
]
% --- protocol: two lifelines, time flows down ---
\node[advbox] (adv) at (-3.7,2.4) {Adversary};
\node[lrnbox] (lrn) at (3.7,2.4) {Learner};
\draw[life] (adv) -- (-3.7,-0.55);
\draw[life] (lrn) -- (3.7,-0.55);
\draw[rstep,thmred!70!black] (-3.7,1.45) -- node[msg]{1.~presents $x_t$} (3.7,1.45);
\draw[rstep,defblue!70!black] (3.7,0.7) -- node[msg]{2.~predicts $\hat y_t$} (-3.7,0.7);
\draw[rstep,thmred!70!black] (-3.7,-0.05) -- node[msg]{3.~reveals $y_t$; mistake if $\hat y_t\neq y_t$} (3.7,-0.05);
\node[font=\scriptsize\itshape, text=black!65] at (0,-0.95) {realizability: $\exists\,h^{\ast}\!\in\mathcal{H},\ h^{\ast}(x_t)=y_t$ for all $t$};
\node[hdl, text=gray!60, anchor=west] at (-6.4,0.7) {round\\$t=1,2,\dots$};
% --- link: why the value is combinatorial ---
\draw[gray!55, ->, densely dashed] (0,-1.3) -- node[hdl, right, pos=0.5]
  {optimal adversary $=$ shattered mistake tree (\cref{fig:mistake-tree})} (0,-2.05);
% --- value chain (what the game is worth) ---
\node[val, fill=sepgreen!10] (opt)  at (-5.2,-2.7) {game value\\$\mathrm{Opt}(\mathcal{H})$};
\node[val, fill=sepgreen!10] (ldim) at (-1.7,-2.7) {$\Ldim(\mathcal{H})$};
\node[val, fill=defblue!8]  (onl)  at (1.8,-2.7) {online\\learnable};
\node[val, fill=white, densely dashed] (pac) at (5.3,-2.7) {PAC\\learnable};
\draw[vedge] (opt) -- node[hdl, above]{$=$} (ldim);
\draw[vedge] (ldim) -- node[hdl, above]{finite $\Leftrightarrow$} (onl);
\draw[sedge] (onl) -- node[hdl, above, text=edgeorange!80!black]{$\Rightarrow$ strict} (pac);
\node[hdl] at (-3.45,-3.45) {\leanname{optimal_mistake_bound_eq_ldim}};
\node[hdl] at (0.05,-3.45) {\leanname{littlestone_characterization}};
\node[hdl, text=edgeorange!75!black] at (3.55,-3.45) {\leanname{online_strictly_stronger_pac}};
\end{tikzpicture}
\caption[The online learning game, whose value is the Littlestone dimension]{The online learning game: round by round the adversary presents an instance, the learner predicts, and the adversary reveals the truth, recording a mistake on every disagreement, with realizability as the only constraint on the adversary. The lower chain shows what the game is worth: the optimal mistake bound equals the Littlestone dimension (\leanname{optimal_mistake_bound_eq_ldim}), the class is online learnable if and only if that dimension is finite, and the one-way implication to PAC learnability is strict, with thresholds on $\R$ (VC dimension one, Littlestone dimension infinite) as the separating witness.}
\label{fig:online-protocol}
\end{figure}

Three features separate this game from the PAC framework of \Cref{ch:pac}, and each one tightens the screws:

\begin{enumerate}[leftmargin=*, itemsep=3pt]
\item \textbf{No distribution.} PAC learning grades the learner against a fixed, unknown distribution $D$; averaging over randomness is baked in. Here there is no distribution. The adversary's choices need not come from any probability measure.

\item \textbf{The adversary adapts.} The adversary watches the learner's predictions and targets the next instance accordingly. This is strictly harder than random sampling: an adaptive adversary can simulate any distribution and then do worse, exploiting whatever the learner just revealed.

\item \textbf{Worst-case over every sequence.} The mistake bound must hold for every sequence the adversary might produce, subject only to realizability. Averaging is off the table.
\end{enumerate}

\begin{defn}[Mistake Bound]
\label{def:mistake-bound}
A learner $A$ is \emph{mistake-bounded} with bound $M$ on $\mathcal{H}$ if, for every adversary strategy consistent with some $h^* \in \mathcal{H}$, the total number of mistakes made by $A$ is at most $M$. The \emph{optimal mistake bound} of $\mathcal{H}$ is $\mathrm{Opt}(\mathcal{H}) = \min_A \max_{\mathrm{adversary}} \#\{\text{mistakes of } A\}$.
\end{defn}

The central question is: what pins $\mathrm{Opt}(\mathcal{H})$? The answer is a combinatorial object, the Littlestone dimension, and it is a tree, not a set.

% ================================================================
% SECTION 2: MISTAKE TREES AND THE LITTLESTONE DIMENSION
% The key combinatorial object. TikZ diagram essential.
% ================================================================

\section{Mistake Trees and the Littlestone Dimension}
\label{sec:mistake-trees}

The VC dimension asks: how large a \emph{set} can $\mathcal{H}$ shatter? The Littlestone dimension asks a harder question: how deep a \emph{tree} can $\mathcal{H}$ shatter? The difference is adaptivity. A shattered set fixes its points before asking whether the class can realize all labelings. A shattered tree lets the points at each depth depend on the labels above, which is exactly the power an adaptive adversary exercises. Replacing sets with trees is the move that captures the right object.

\begin{defn}[Mistake Tree]
\label{def:mistake-tree}
A \emph{mistake tree} for $\mathcal{H}$ over $X$ is a complete binary tree $T$ in which:
\begin{itemize}[leftmargin=*, itemsep=2pt]
\item Each internal node $v$ is labeled with an instance $x_v \in X$.
\item The left child of $v$ corresponds to the label $y = 0$ and the right child to $y = 1$.
\item For every root-to-leaf path $(v_1, y_1), (v_2, y_2), \ldots, (v_d, y_d)$, there exists a hypothesis $h \in \mathcal{H}$ consistent with all labels along the path: $h(x_{v_i}) = y_i$ for all $i$.
\end{itemize}
The \emph{depth} of the tree is the number of internal nodes on any root-to-leaf path (all paths have the same length since the tree is complete).\formal{FLT_Proofs.Complexity.GameInfra.LTree.isShattered}
\end{defn}

A mistake tree is precisely an adversary strategy written down as a tree. Starting at the root, the adversary presents $x_{v_1}$. Whatever the learner predicts, the adversary labels $x_{v_1}$ to contradict that prediction and descends to the corresponding child. Because every root-to-leaf path is consistent with some $h \in \mathcal{H}$, realizability holds throughout. A tree shattered at depth $d$ therefore witnesses an adversary strategy that forces \emph{any} learner to make at least $d$ mistakes.

\begin{defn}[Littlestone Dimension {\cite{Littlestone1988}}]
\label{def:ldim}
The \emph{Littlestone dimension} of a hypothesis class $\mathcal{H}$, denoted $\Ldim(\mathcal{H})$, is the maximum depth of a complete mistake tree for $\mathcal{H}$. If arbitrarily deep mistake trees exist, $\Ldim(\mathcal{H}) = \infty$.\formal{FLT_Proofs.Complexity.GameInfra.LittlestoneDim}
\end{defn}

\begin{figure}[ht]
\centering
\begin{tikzpicture}[
    level distance=1.6cm,
    sibling distance=1.2cm,
    every node/.style={font=\small},
    treenode/.style={draw, circle, thick, minimum size=0.7cm, fill=onlineblue!10, font=\small\bfseries},
    leafnode/.style={draw, rectangle, rounded corners, thick, minimum width=0.6cm, minimum height=0.45cm, fill=sepgreen!15, font=\scriptsize},
    edgelabel/.style={font=\scriptsize, midway},
    advedge/.style={draw=thmred, line width=1.3pt},
    level 1/.style={sibling distance=5.6cm},
    level 2/.style={sibling distance=2.8cm},
    level 3/.style={sibling distance=1.4cm},
]
\node[treenode] (root) {$4$}
    child { node[treenode] (n2) {$2$}
        child { node[treenode] (n1) {$1$}
            child { node[leafnode] (h0) {$h_{\leq 0}$} edge from parent node[edgelabel, left] {$0$} }
            child { node[leafnode] {$h_{\leq 1}$} edge from parent node[edgelabel, right] {$1$} }
            edge from parent node[edgelabel, left] {$0$}
        }
        child { node[treenode] {$3$}
            child { node[leafnode] {$h_{\leq 2}$} edge from parent node[edgelabel, left] {$0$} }
            child { node[leafnode] {$h_{\leq 3}$} edge from parent node[edgelabel, right] {$1$} }
            edge from parent node[edgelabel, right] {$1$}
        }
        edge from parent node[edgelabel, left] {$0$}
    }
    child { node[treenode] {$6$}
        child { node[treenode] {$5$}
            child { node[leafnode] {$h_{\leq 4}$} edge from parent node[edgelabel, left] {$0$} }
            child { node[leafnode] {$h_{\leq 5}$} edge from parent node[edgelabel, right] {$1$} }
            edge from parent node[edgelabel, left] {$0$}
        }
        child { node[treenode] {$7$}
            child { node[leafnode] {$h_{\leq 6}$} edge from parent node[edgelabel, left] {$0$} }
            child { node[leafnode] {$h_{\leq 7}$} edge from parent node[edgelabel, right] {$1$} }
            edge from parent node[edgelabel, right] {$1$}
        }
        edge from parent node[edgelabel, right] {$1$}
    };
% --- highlighted adversary run: root -> 2 -> 1 -> h_{<=0}, forces depth-many (3) mistakes ---
\draw[advedge] (root) -- (n2);
\draw[advedge] (n2) -- (n1);
\draw[advedge] (n1) -- (h0);
\node[font=\scriptsize\itshape, text=thmred, anchor=north, align=center] at ($(h0.south)+(0,-0.12)$)
  {adversary run:\\$3$ mistakes forced};
% --- the general meaning + verified handles (A4: verified >= vs borrowed =) ---
\node[font=\scriptsize, align=left, anchor=north west] at (3.7,1.05) {%
  \textbf{A shattered mistake tree.}\\[2pt]
  \textit{$X=\{1,\dots,8\}$; thresholds $h_{\leq\theta}(x)=\ind[x\leq\theta]$.}\\[1pt]
  \textit{internal node: an instance; leaf: a consistent $h$.}\\[1pt]
  \textit{every path realizable $\Rightarrow$ tree shattered.}\\[3pt]
  depth $=3=\Ldim(\mathcal{H})$\quad\leanname{LittlestoneDim}\\[1pt]
  a depth-$d$ shattered tree forces $d$ mistakes:\\
  \quad\leanname{adversary_core}\\[2pt]
  here $\Ldim\geq 3$\quad\leanname{thresholdTree_shattered}\\[1pt]
  \textcolor{notegray}{$=\lceil\log_2 8\rceil$\quad(\cite{Littlestone1988})}
};
\end{tikzpicture}
\caption[A shattered mistake tree of depth three witnessing $\Ldim=3$ for thresholds on eight points]{A complete mistake tree of depth $3$ for thresholds on $\{1,\ldots,8\}$, witnessing $\Ldim(\mathcal{H}) \geq 3$ (\leanname{thresholdTree_shattered}). Each internal node is an instance; each edge is a label; every root-to-leaf path is realized by the threshold named at its leaf. The highlighted run shows the adversary executing binary search: at each node it assigns the label opposite to the learner's prediction and descends, forcing $3$ mistakes in $3$ rounds while maintaining realizability throughout (\leanname{adversary_core}). The Littlestone dimension is the maximum depth of any such shattered tree (\leanname{LittlestoneDim}); the fact that this tree has depth $\lceil\log_2 8\rceil = 3$ confirms the threshold bound.}
\label{fig:mistake-tree}
\end{figure}

\begin{example}[Littlestone dimension of fundamental classes]
\label{ex:ldim-examples}
\begin{enumerate}[label=(\alph*), leftmargin=*, itemsep=4pt]
\item \textbf{Finite classes.} If $|\mathcal{H}| < \infty$, then $\Ldim(\mathcal{H}) \leq \lfloor \log_2 |\mathcal{H}| \rfloor$, since a complete binary tree of depth $d$ has $2^d$ leaves and each leaf requires a distinct consistent hypothesis. This bound is tight for classes that are ``richly structured enough'' (e.g., all functions on a domain of size $\log_2 |\mathcal{H}|$).

\item \textbf{Thresholds on $\{1, \ldots, n\}$.} The class $\{h_{\leq \theta} : \theta \in \{0, 1, \ldots, n\}\}$ has $\Ldim = \lceil \log_2(n+1) \rceil$. The mistake tree in \Cref{fig:mistake-tree} illustrates the case $n = 8$. The adversary performs binary search on the threshold.

\item \textbf{Thresholds on $\R$.} Now $\Ldim = \infty$. The key point: the domain is dense, so the adversary can always find a point between any two previous thresholds. At each round, the adversary presents a point that bisects the current interval of uncertainty, forcing a mistake no matter what the learner predicts. This produces mistake trees of unbounded depth.

\item \textbf{Intervals on $\R$.} The class $\{x \mapsto \ind[a \leq x \leq b] : a, b \in \R\}$ also has $\Ldim = \infty$. The adversary can adaptively narrow down both endpoints.

\item \textbf{Linear classifiers in $\R^d$.} The class of halfspaces has $\Ldim = d$, matching the VC dimension. This is one of the rare cases where the two dimensions coincide.
\end{enumerate}
\end{example}

\begin{remark}[Sets vs.\ trees]
\label{rmk:sets-vs-trees}
Shattering a \emph{set} $\{x_1, \ldots, x_d\}$ means $\mathcal{H}$ can produce all $2^d$ labelings of these \emph{fixed} points. Shattering a \emph{tree} of depth $d$ means $\mathcal{H}$ can produce a consistent labeling along every root-to-leaf path, but the points themselves may \emph{depend on previous labels}. The tree structure captures adaptivity: the adversary's choice at depth $k$ depends on the learner's responses at depths $1, \ldots, k-1$. A shattered set is a special case of a mistake tree (one in which every node at the same depth is labeled with the same instance), so $\VC(\mathcal{H}) \leq \Ldim(\mathcal{H})$ always holds. The converse fails: thresholds on $\R$ have $\VC = 1$ but $\Ldim = \infty$.
\end{remark}

% ================================================================
% SECTION 3: THE STANDARD OPTIMAL ALGORITHM (SOA)
% The constructive heart of online learning. Read like describing an algorithm.
% ================================================================

\section{The Standard Optimal Algorithm}
\label{sec:soa}

The simplest reasonable strategy is Halving: maintain the version space (all hypotheses consistent with everything seen so far), predict the majority vote. Every mistake eliminates at least half the surviving hypotheses, so Halving makes at most $\lfloor \log_2 |\mathcal{H}| \rfloor$ mistakes when $\mathcal{H}$ is finite.

For infinite classes, cardinality collapses and Halving offers nothing. The Standard Optimal Algorithm (SOA), due to Littlestone~\cite{Littlestone1988}, replaces ``majority vote'' with a sharper criterion: predict the label whose sub-version-space carries the larger \emph{Littlestone dimension}. The switch from counting hypotheses to measuring mistake-tree depth is what makes the algorithm optimal for every hypothesis class, finite or infinite.

\begin{defn}[Standard Optimal Algorithm (\SOA)]
\label{def:soa}
Given a hypothesis class $\mathcal{H}$, the $\SOA$ maintains a version space $V_t \subseteq \mathcal{H}$, initially $V_1 = \mathcal{H}$. At round $t$, upon receiving instance $x_t$:
\begin{enumerate}[leftmargin=*, itemsep=2pt]
\item Partition $V_t$ into $V_t^0 = \{h \in V_t : h(x_t) = 0\}$ and $V_t^1 = \{h \in V_t : h(x_t) = 1\}$.
\item Predict $\hat{y}_t = \arg\max_{b \in \{0,1\}} \Ldim(V_t^b)$.
\item After observing the true label $y_t$, update $V_{t+1} = V_t^{y_t}$.
\end{enumerate}
Ties in step 2 are broken arbitrarily.\formal{FLT_Proofs.Complexity.GameInfra.SOA}
\end{defn}

The SOA reasons about future adversarial damage. At each round it asks: if this prediction is wrong, which sub-version-space does the adversary inherit, and how deeply can they search from there? By predicting the label whose sub-version-space has the larger Littlestone dimension, the algorithm guarantees that every mistake lands it in the smaller sub-version-space, reducing the Littlestone dimension by at least one. The game ends when the dimension reaches zero.

\begin{thm}[Upper Bound: $\SOA$ achieves $\Ldim(\mathcal{H})$ mistakes {\cite{Littlestone1988}}]
\label{thm:soa-upper}
For any hypothesis class $\mathcal{H}$ with $\Ldim(\mathcal{H}) = d < \infty$, the $\SOA$ makes at most $d$ mistakes against any adversary.\formal{FLT_Proofs.Theorem.Online.soa_mistakes_bounded}
\end{thm}

\begin{proof}
We show that the Littlestone dimension of the version space drops by at least one at each mistake. Define $d_t = \Ldim(V_t)$. We claim:
\begin{equation}
\label{eq:ldim-drop}
\text{If the $\SOA$ makes a mistake at round } t, \text{ then } d_{t+1} \leq d_t - 1.
\end{equation}%
\formal{FLT_Proofs.Theorem.Online.ldim_strict_decrease_on_mistake}

\emph{Proof of the claim.} Suppose the $\SOA$ makes a mistake at round $t$. This means $\hat{y}_t \neq y_t$. By the algorithm's rule, $\hat{y}_t$ was chosen so that $\Ldim(V_t^{\hat{y}_t}) \geq \Ldim(V_t^{y_t})$, i.e., the $\SOA$ predicted the side with the \emph{larger} (or equal) Littlestone dimension. After the mistake, the version space updates to $V_{t+1} = V_t^{y_t}$, the side with the \emph{smaller} (or equal) Littlestone dimension.

Now we must show that $\Ldim(V_t^{y_t}) \leq d_t - 1$. Suppose for contradiction that $\Ldim(V_t^0) \geq d_t$ and $\Ldim(V_t^1) \geq d_t$. Then there exists a complete mistake tree $T_0$ of depth $d_t$ for $V_t^0$ and a complete mistake tree $T_1$ of depth $d_t$ for $V_t^1$. We can construct a complete mistake tree of depth $d_t + 1$ for $V_t$: the root is labeled $x_t$, its left subtree (corresponding to label $0$) is $T_0$, and its right subtree (corresponding to label $1$) is $T_1$. Every root-to-leaf path through the left subtree is consistent with some $h \in V_t^0 \subseteq V_t$ (and $h(x_t) = 0$), and similarly for the right subtree. This yields $\Ldim(V_t) \geq d_t + 1$, contradicting $\Ldim(V_t) = d_t$.

Therefore $\min\{\Ldim(V_t^0), \Ldim(V_t^1)\} \leq d_t - 1$. Since $\hat{y}_t$ selects the side with the larger dimension, $V_{t+1} = V_t^{y_t}$ is the side with the smaller dimension, so $d_{t+1} = \Ldim(V_{t+1}) \leq d_t - 1$.

\emph{Completing the proof.} We have $d_1 = \Ldim(\mathcal{H}) = d$. Each mistake decreases $d_t$ by at least $1$. Since the Littlestone dimension is a non-negative integer (the version space always contains the target $h^*$, so $V_t \neq \emptyset$ and $\Ldim(V_t) \geq 0$), the total number of mistakes is at most $d$.
\end{proof}

\begin{remark}[SOA vs.\ Halving]
\label{rmk:soa-halving}
For finite $\mathcal{H}$, the Halving algorithm predicts by majority vote: $\hat{y}_t = \arg\max_b |V_t^b|$. Each mistake halves the version space, giving at most $\lfloor \log_2 |\mathcal{H}| \rfloor$ mistakes. The $\SOA$ replaces cardinality $|V_t^b|$ with $\Ldim(V_t^b)$. For finite classes, this distinction is often unimportant (both give $O(\log |\mathcal{H}|)$ mistakes), but for infinite classes, cardinality is meaningless while the Littlestone dimension is finite and well-behaved.
\end{remark}

\begin{remark}[Computability of the $\SOA$]
\label{rmk:soa-computability}
The $\SOA$ is an \emph{information-theoretically} optimal algorithm, but it is not necessarily \emph{computationally} efficient. Computing $\Ldim(V_t^b)$ at each round may be undecidable for some hypothesis classes. The $\SOA$ is thus an existence proof: it shows that $d$ mistakes suffice, even if finding the optimal prediction at each step is computationally hard. Efficient online algorithms for specific classes (Perceptron for halfspaces, Winnow for disjunctions) achieve near-optimal mistake bounds through class-specific structure.
\end{remark}

% ================================================================
% SECTION 4: THE LOWER BOUND - ADVERSARY STRATEGY
% The adversary plays the mistake tree from root to leaf.
% ================================================================

\section{The Lower Bound: The Adversary Strategy}
\label{sec:lower-bound}

The SOA limits the learner's mistakes to $\Ldim(\mathcal{H})$. That is the ceiling. The lower bound closes the gap: no learner can stay below that ceiling, because for any learner an adversary exists that forces at least $\Ldim(\mathcal{H})$ mistakes. Together the two bounds pin $\mathrm{Opt}(\mathcal{H})$ exactly.

\begin{thm}[Lower Bound: Any learner makes at least $\Ldim(\mathcal{H})$ mistakes {\cite{Littlestone1988}}]
\label{thm:lower-bound}
For any hypothesis class $\mathcal{H}$ with $\Ldim(\mathcal{H}) = d$ and any (possibly randomized) learner $A$, there exists an adversary strategy that forces $A$ to make at least $d$ mistakes.\formal{FLT_Proofs.Theorem.Online.adversary_lower_bound}
\end{thm}

\begin{proof}
Let $T$ be a complete mistake tree for $\mathcal{H}$ of depth $d$. The adversary plays $T$ as follows.

The adversary maintains a pointer to a node of $T$, starting at the root. At round $t$, the adversary presents the instance $x_{v_t}$ labeling the current node $v_t$. The learner predicts $\hat{y}_t$. The adversary then sets $y_t = 1 - \hat{y}_t$ (the opposite of the learner's prediction) and descends to the child of $v_t$ corresponding to label $y_t$.

This strategy has two properties:
\begin{enumerate}[leftmargin=*, itemsep=2pt]
\item \textbf{Every round is a mistake.} By construction, $y_t = 1 - \hat{y}_t \neq \hat{y}_t$.
\item \textbf{Realizability is maintained.} After $d$ rounds, the adversary has descended from the root to a leaf of $T$, tracing a path $(v_1, y_1), (v_2, y_2), \ldots, (v_d, y_d)$. By the definition of a mistake tree, there exists $h^* \in \mathcal{H}$ consistent with all labels along this path: $h^*(x_{v_i}) = y_i$ for all $i \in \{1, \ldots, d\}$. For any subsequent rounds $t > d$, the adversary can label consistently with this $h^*$.
\end{enumerate}

The adversary forces exactly $d$ mistakes in the first $d$ rounds, so $A$ makes at least $d$ mistakes.

For randomized learners, the argument extends by the minimax theorem (or directly): for any distribution over predictions $\hat{y}_t$, the adversary's strategy of playing the opposite forces an expected mistake at every round. Alternatively, one can apply Yao's minimax principle: the worst-case deterministic adversary against the best randomized learner equals the best deterministic learner against the worst-case adversary, and we have shown the latter is at least $d$.
\end{proof}

Combining \Cref{thm:soa-upper,thm:lower-bound}, we obtain the fundamental characterization.

\begin{thm}[Littlestone's Characterization {\cite{Littlestone1988}}]
\label{thm:littlestone-characterization}
For any hypothesis class $\mathcal{H}$:
\[
\mathrm{Opt}(\mathcal{H}) = \Ldim(\mathcal{H}).
\]
That is, the optimal mistake bound of $\mathcal{H}$ in the online learning game is exactly the Littlestone dimension of $\mathcal{H}$. In particular, $\mathcal{H}$ admits a finite mistake bound if and only if $\Ldim(\mathcal{H}) < \infty$.\formal{FLT_Proofs.Theorem.Online.optimal_mistake_bound_eq_ldim}
\end{thm}

\begin{historical}
Littlestone proved this characterization in 1988~\cite{Littlestone1988}, just four years after Valiant's PAC model. The result established online learning as a mathematically independent paradigm: the combinatorial quantity governing online learnability is genuinely different from the one governing PAC learnability. The $\SOA$ is sometimes called the ``Standard Optimal Algorithm'' precisely because it achieves the information-theoretic optimum; the name is due to the online learning community's convention.
\end{historical}

% ================================================================
% SECTION 5: REGRET BOUNDS AND MULTIPLICATIVE WEIGHTS
% Profile C (Revelation). Regret is surprising because it is relative.
% ================================================================

\section{Regret Bounds and the Multiplicative Weights Framework}
\label{sec:regret}

The mistake-bound framework rests on realizability: some $h^* \in \mathcal{H}$ must be consistent with every label. That is a substantial assumption. When it fails, the optimal mistake count is infinite and the framework offers no useful bound.

The \emph{regret} framework drops realizability and changes the benchmark. Rather than asking for an absolute mistake count, it grades the learner against the \emph{best fixed hypothesis in hindsight}: whoever single $h \in \mathcal{H}$ would have made the fewest errors on the realized sequence, even if that $h$ errs constantly. The learner's job is to stay close to that benchmark, not to be right.

\begin{defn}[Regret]
\label{def:regret}
Given a sequence $(x_1, y_1), \ldots, (x_T, y_T)$ and the learner's predictions $\hat{y}_1, \ldots, \hat{y}_T$, the \emph{regret} of the learner relative to $\mathcal{H}$ is
\[
\mathrm{Regret}_T = \sum_{t=1}^T \ind[\hat{y}_t \neq y_t] - \min_{h \in \mathcal{H}} \sum_{t=1}^T \ind[h(x_t) \neq y_t].
\]
\end{defn}

The definition holds a conceptual surprise. Regret can be small even when the learner makes many mistakes, because the benchmark also makes many mistakes. If the adversary is unrealizable and every hypothesis in $\mathcal{H}$ errs frequently, then keeping pace with the best of them is a genuine achievement, even with a large absolute error count. What the regret framework measures is the \emph{cost of not knowing which hypothesis is best}, not the absolute cost of prediction.

\begin{remark}[Mistake bound vs.\ regret bound]
\label{rmk:mistake-vs-regret}
In the realizable case, the best hypothesis $h^*$ makes zero mistakes, so $\mathrm{Regret}_T$ equals the total mistake count. The regret framework is therefore a strict generalization of the mistake-bound framework. The regret formulation becomes essential when $\mathcal{H}$ is infinite or when realizability fails: precisely the settings where the Littlestone dimension may be infinite and the mistake-bound framework provides no guarantee.
\end{remark}

The Weighted Majority algorithm of Littlestone and Warmuth~\cite{Littlestone1988} achieves the following regret bound for finite $\mathcal{H}$.

\begin{thm}[Weighted Majority / Hedge]
\label{thm:weighted-majority}
For a finite hypothesis class $\mathcal{H}$ with $|\mathcal{H}| = N$, the Weighted Majority algorithm achieves, for any sequence of $T$ rounds:
\[
\mathrm{Regret}_T \leq 2\sqrt{T \ln N}.
\]
In particular, the per-round regret $\mathrm{Regret}_T / T \to 0$ as $T \to \infty$.
\end{thm}

The algorithm maintains a weight $w_t(h)$ for each $h \in \mathcal{H}$, initially $w_1(h) = 1$. At each round: predict the weighted majority vote; after observing $y_t$, multiply the weight of each hypothesis that erred by a factor $(1 - \eta)$ for a learning rate $\eta \in (0, 1)$. Setting $\eta = \sqrt{\ln N / T}$ (or using the doubling trick when $T$ is unknown) yields the bound above.

This approach generalizes to the \emph{multiplicative weights} framework (also called Hedge), which extends beyond binary prediction to decision-making over finitely many ``experts.'' The framework is one of the most widely applicable algorithmic paradigms in theoretical computer science, with applications to zero-sum games, boosting, and linear programming~\cite{ShalevShwartzBenDavid2014}.

\begin{remark}[Connection to Littlestone dimension]
\label{rmk:regret-ldim}
Ben-David, P\'al, and Shalev-Shwartz~\cite{BenDavidPalShalevShwartz2009} showed that for infinite hypothesis classes, the Littlestone dimension also governs the optimal regret rate. Specifically, the minimax regret over $T$ rounds is $\Theta(\sqrt{\Ldim(\mathcal{H}) \cdot T})$ in the agnostic (non-realizable) setting. This connects the combinatorial tree structure of the Littlestone dimension to the sequential prediction framework even beyond realizability.
\end{remark}

% ================================================================
% SECTION 6: THE PAC-ONLINE GAP
% VC ≤ Ldim always; thresholds show the gap can be infinite.
% ================================================================

\subsection{Attention as exponential weights}
\label{sec:attention-online}

The multiplicative-weights update of \cref{thm:weighted-majority} is the computational core of the attention mechanism.  Fix scores $s(x) \in \R^k$ assigned to $k$ keys by a parametric router on input $x$.  The soft-attention weight on key $i$ is the softmax $w_i(x) = e^{s_i(x)} / \sum_j e^{s_j(x)}$, which is exactly the exponential-weights (Hedge) distribution with the scores in the role of cumulative gains.

\begin{prop}[Softmax attention is an exponential-weights step]
\label{prop:attention-hedge}
The soft-attention weights $w_i(x) \propto \exp(s_i(x))$ are the Hedge distribution of \cref{thm:weighted-majority} over the $k$ keys.\formal{TLT_Proofs.Routing.MeasurableScoreRouting.softmaxWeight} The weight order matches the score order ($w_i(x) \le w_j(x)$ if and only if $s_i(x) \le s_j(x)$), so the hard, argmax route attends to the maximally-weighted key, the Halving-style greedy specialization of the same update.\formal{TLT_Proofs.Routing.MeasurableScoreRouting.softmaxWeight_le_iff_score}
\end{prop}

\begin{proof}
Both claims unfold the softmax.  The normalization $\sum_j e^{s_j(x)}$ is positive, so $w_i \propto e^{s_i}$ is a probability vector, and it coincides termwise with the exponential-weights distribution that assigns mass $\propto e^{\eta G_i}$ to expert $i$ of cumulative gain $G_i$ at learning rate $\eta$ (the inverse temperature, absorbed into the score scale).  Since $t \mapsto e^t$ is strictly increasing, $w_i \le w_j \Leftrightarrow e^{s_i} \le e^{s_j} \Leftrightarrow s_i \le s_j$, and the $\arg\max$ of the weights is the $\arg\max$ of the scores.
\end{proof}

The correspondence makes the regret framework of \cref{sec:regret} the right setting for a routed transformer: a layer attending over $k$ keys runs one exponential-weights step, and the regret bound $2\sqrt{T \ln k}$ bounds how much worse the routed predictions are than the single best key in hindsight.

\begin{remark}[Route stability]
\label{rmk:route-stability}
A tempered attention router's hard route is stable: away from a score tie, a small perturbation of the scores, or of the parameters and temperature that produce them, leaves the routed key unchanged.\formal{TLT_Proofs.TemperedDesignLaw.Stability.TD7_route_stable_proof} The inputs at which the route flips form the score-tie boundary, of the same measure-zero character that the mistake-tree analysis of \cref{sec:soa} exploits when it splits the version space by a single instance's label.  Stability is what lets the discrete, argmax-routed learner inherit the regret guarantee of its smooth softmax surrogate.
\end{remark}

\section{The PAC--Online Gap}
\label{sec:pac-online-gap}

Both the VC dimension and the Littlestone dimension measure the complexity of a hypothesis class, but each captures a different aspect. The VC dimension asks what the class can do to a fixed set of points; the Littlestone dimension asks what an adaptive adversary can do when the points are chosen sequentially in response to the learner. The two measures agree on some classes and diverge wildly on others.

\begin{prop}[$\VC \leq \Ldim$]
\label{prop:vc-leq-ldim}
For any hypothesis class $\mathcal{H}$, $\VC(\mathcal{H}) \leq \Ldim(\mathcal{H})$.\formal{FLT_Proofs.Complexity.Littlestone.BranchWiseLittlestoneDim_ge_VCDim}
\end{prop}

\begin{proof}
Let $S = \{x_1, \ldots, x_d\}$ be a set shattered by $\mathcal{H}$, so $\VC(\mathcal{H}) \geq d$. We construct a complete mistake tree of depth $d$. The root is labeled $x_1$. Every node at depth $k$ is labeled $x_{k+1}$ (the same instance at every node of a given depth). For any root-to-leaf path with labels $(y_1, \ldots, y_d)$, shattering guarantees that some $h \in \mathcal{H}$ satisfies $h(x_i) = y_i$ for all $i$. Thus this is a valid mistake tree of depth $d$, so $\Ldim(\mathcal{H}) \geq d$.
\end{proof}

\begin{figure}[ht]
\centering
\begin{tikzpicture}[
    level distance=1.4cm,
    every node/.style={font=\small},
    treenode/.style={draw, circle, thick, minimum size=0.65cm, fill=defblue!10, font=\small},
    leafnode/.style={draw, rectangle, rounded corners, thick, minimum width=0.5cm, minimum height=0.4cm, fill=sepgreen!15, font=\scriptsize},
    edgelabel/.style={font=\scriptsize, midway},
    setbox/.style={draw, thick, rounded corners, fill=defblue!6, align=center, font=\scriptsize, inner sep=4pt},
    embed/.style={-{Stealth[length=2.5mm]}, thick, defblue!75!black},
    gap/.style={-{Stealth[length=2.5mm]}, thick, densely dashed, thmred!80!black},
    hdl/.style={font=\tiny, text=black!70, align=center},
    level 1/.style={sibling distance=3.6cm},
    level 2/.style={sibling distance=1.8cm},
]
% --- the non-adaptive tree: same instance x_2 at every depth-1 node ---
\node[treenode] (root) {$x_1$}
    child { node[treenode] {$x_2$}
        child { node[leafnode] {$(0,0)$} edge from parent node[edgelabel, left] {$0$} }
        child { node[leafnode] {$(0,1)$} edge from parent node[edgelabel, right] {$1$} }
        edge from parent node[edgelabel, left] {$0$}
    }
    child { node[treenode] {$x_2$}
        child { node[leafnode] {$(1,0)$} edge from parent node[edgelabel, left] {$0$} }
        child { node[leafnode] {$(1,1)$} edge from parent node[edgelabel, right] {$1$} }
        edge from parent node[edgelabel, right] {$1$}
    };
% --- (A) shattered set source + solid embedding arrow (set -> non-adaptive tree) ---
\node[setbox] (set) at (-4.8,-0.35) {Shattered set\\$S=\{x_1,x_2\}$\\\textsc{vc}: all $2^2$\\labelings realized};
\node[font=\scriptsize\itshape, text=notegray] at (-1.95,0.95) {embed (same instance per depth)};
\node[font=\small] at (-1.95,0.55) {$\VC(\mathcal{H})\le\Ldim(\mathcal{H})$};
\node[hdl] at (-1.95,0.18) {\leanname{BranchWiseLittlestoneDim_ge_VCDim}};
\draw[embed] (set.east) to[out=0,in=168] (root.west);
\node[hdl, text=notegray, align=center, anchor=north] at (-4.8,-1.3) {(branch-wise $=$ path-\\consistent on this\\non-adaptive tree)};
% --- (C) strict-gap dashed arrow + note: adaptivity escapes the image ---
\node[draw, thick, densely dashed, thmred!75!black, rounded corners, fill=thmred!4,
      align=left, font=\scriptsize, inner sep=4pt] (gapnote) at (5.0,-2.3)
  {\textbf{One-way only.}\\an adaptive tree need not\\come from a set: thresholds\\on $\R$ give $\VC=1$, $\Ldim=\infty$.\\[2pt]
   {\tiny\leanname{ldim_threshold_top}}\\{\tiny\leanname{online_strictly_stronger_pac}}};
\draw[gap] (2.0,-1.7) to[out=-12,in=165] (gapnote.west);
% --- (D) one-liner under the tree ---
\node[font=\scriptsize\itshape, text=notegray, align=center, anchor=north, text width=10.5cm] at (-0.3,-3.5)
  {Non-adaptive: the same instance at each depth, so the tree \emph{is} the set; dropping that constraint is adaptivity, which is what lets $\Ldim$ exceed $\VC$.};
\end{tikzpicture}
\caption[Shattered set embeds as a non-adaptive tree, proving $\VC\le\Ldim$; thresholds on $\R$ show the gap is infinite]{The argument for $\VC(\mathcal{H})\le\Ldim(\mathcal{H})$ made visible: a shattered set $\{x_1,x_2\}$ becomes a non-adaptive mistake tree of the same depth by placing the same instance at every node of each level (\leanname{BranchWiseLittlestoneDim_ge_VCDim}). Every root-to-leaf path is realized because shattering supplies a consistent hypothesis for each labeling. The embedding is one-way: an adaptive tree lets the instances vary with the labels, escaping the image of any shattered set. Thresholds on $\R$ demonstrate that the gap is unbounded, with $\VC=1$ and $\Ldim=\infty$ (\leanname{ldim_threshold_top}, \leanname{online_strictly_stronger_pac}). The slack $\Ldim - \VC$ is the precise measure of what adaptivity adds.}
\label{fig:vc-leq-ldim}
\end{figure}

The inequality $\VC \leq \Ldim$ is tight for some classes: for halfspaces in $\R^d$, both dimensions equal $d$. But the gap can be infinite, and a single class witnesses it.

\begin{prop}[The PAC--Online separation: thresholds on $\R$]
\label{prop:pac-online-sep}
Let $\mathcal{H}_{\mathrm{thr}} = \{x \mapsto \ind[x \leq \theta] : \theta \in \R\}$. Then $\VC(\mathcal{H}_{\mathrm{thr}}) = 1$ and $\Ldim(\mathcal{H}_{\mathrm{thr}}) = \infty$.\formal{FLT_Proofs.Theorem.Separation.pac_not_implies_online}
\end{prop}

\begin{proof}
$\VC = 1$: Any single point $x$ is shattered (take thresholds $\theta < x$ and $\theta > x$). No two points $x_1 < x_2$ are shattered, because $h(x_1) = 0, h(x_2) = 1$ requires $x_2 \leq \theta < x_1$, which contradicts $x_1 < x_2$.

$\Ldim = \infty$: We construct mistake trees of arbitrary depth. For any $d$, consider the domain restricted to $\{1, 2, \ldots, 2^d\}$. The adversary can perform binary search: at depth $k$, present the midpoint of the current interval. Whatever the learner predicts, the adversary labels to force a mistake (and narrows the interval by half). After $d$ rounds, the adversary has traced a path from root to leaf, consistent with the threshold at the boundary of the final interval. Since $d$ was arbitrary, $\Ldim = \infty$.
\end{proof}

This separation is the most important non-implication in the concept graph. It witnesses the edge $\nde{pac\_learning} \xrightarrow{\rel{does\_not\_imply}} \nde{online\_learning}$: a class can be PAC learnable (finite VC dimension) yet not online learnable (infinite Littlestone dimension). The converse implication does hold: $\Ldim < \infty$ implies $\VC < \infty$ (by $\VC \leq \Ldim$), so online learnability implies PAC learnability.\formal{FLT_Proofs.Theorem.Separation.online_imp_pac}

\begin{separation}
\textbf{PAC $\not\Rightarrow$ Online.} Witness: $\mathcal{H}_{\mathrm{thr}}$ on $\R$. In the online game, the adversary uses the order structure of $\R$ to perform binary search on the threshold, forcing a mistake at every round. A random sample from a fixed distribution cannot exploit this structure; with high probability, the sample locates the threshold to within $\varepsilon$. The gap between PAC and online learnability reflects the difference between a world that is fixed and a world that watches you.
\end{separation}

\begin{remark}[Why the gap arises]
\label{rmk:gap-intuition}
In PAC learning the distribution $D$ is fixed before the game begins; the learner faces a static environment. In online learning the adversary adapts to the learner's behavior, and adaptivity is the source of the gap. A threshold on $\R$ is easy to learn from random data: one sample locates it approximately. Against an adaptive adversary the same class is impossible to learn: the adversary always presents a point that bisects the remaining uncertainty, and the bisection never ends. The Littlestone dimension measures the depth of that adaptive search, and it can exceed the VC dimension by an unbounded amount. A full treatment of this and related separations appears in \Cref{ch:separations}.
\end{remark}

\section{What This Chapter Established}
\label{sec:ch6-summary}

The sequential game between learner and adversary pins its value to a single combinatorial quantity:

\begin{enumerate}[leftmargin=*, itemsep=3pt]
\item \textbf{The Littlestone characterization.} The optimal mistake bound for online learning equals the Littlestone dimension $\Ldim(\mathcal{H})$, the maximum depth of a complete mistake tree. This is a \emph{combinatorial} characterization, in contrast to the VC dimension's \emph{set-combinatorial} characterization of PAC learning.

\item \textbf{The SOA algorithm.} The upper bound is constructive: the Standard Optimal Algorithm achieves $\Ldim(\mathcal{H})$ mistakes by predicting the label whose version space has the larger Littlestone dimension. Each mistake provably reduces the Littlestone dimension of the version space by at least one.

\item \textbf{The adversary strategy.} The lower bound is also constructive: the adversary traverses a complete mistake tree from root to leaf, labeling each instance to contradict the learner's prediction.

\item \textbf{The PAC--Online gap.} $\VC(\mathcal{H}) \leq \Ldim(\mathcal{H})$ always, but the gap can be infinite (thresholds: $\VC = 1$, $\Ldim = \infty$). Online learnability implies PAC learnability, but not conversely.

\item \textbf{The regret framework.} Regret bounds, measuring performance relative to the best fixed hypothesis, extend online learning beyond the realizable setting. The Littlestone dimension governs optimal regret rates even in the agnostic case.
\end{enumerate}

The combinatorial and algorithmic character of these results is structural. Online learning theory runs on trees and explicit strategies; PAC learning theory runs on sets and tail probabilities. Both measure the complexity of the same object, a hypothesis class $\mathcal{H}$, but the instruments are genuinely different: the adversary adapts to the learner's behavior round by round, while a fixed distribution cannot.

\section{Open Problems}
\label{sec:ch6-open}

Each problem is anchored to a result above and tagged: a \emph{known unknown} (\textsc{ku}) is an articulated open question; an \emph{unknown known} (\textsc{uk}) is a result in the literature not yet absorbed into a clean statement.

\begin{openprob}[The Littlestone dimension of attention, \textsc{uk}]
\label{op:ldim-attention}
\Cref{prop:attention-hedge} identifies soft attention with exponential weights.  Determine the Littlestone dimension of a finite-head attention class as a function of the architecture and the inverse temperature: is it finite (so the class is online learnable, \cref{thm:littlestone-characterization}), and does it diverge as the temperature grows and the router approaches a hard threshold, the regime in which thresholds on $\R$ already give $\Ldim = \infty$ (\cref{prop:pac-online-sep})?
\end{openprob}

\begin{openprob}[Regret of online attention routing, \textsc{uk}]
\label{op:attention-regret}
Via \cref{prop:attention-hedge}, a layer routing over $k$ keys runs one Hedge step.  Prove the $O(\sqrt{T \ln k})$ regret bound for online attention routing with the inverse temperature as the learning rate, and determine the optimal temperature schedule, connecting it to the route-stability margin (\cref{rmk:route-stability}, \leanname{TLT_Proofs.TemperedDesignLaw.Stability.TD7_route_stable_proof}).
\end{openprob}

\begin{openprob}[Computable optimal online learners, \textsc{ku}]
\label{op:soa-computable}
\Cref{rmk:soa-computability} notes the $\SOA$ may be uncomputable because $\Ldim(V_t^b)$ can be undecidable.  Characterize the classes for which the $\SOA$ is computable in polynomial time per round, and decide whether every class of finite Littlestone dimension admits \emph{some} efficient optimal online learner, or whether an information--computation gap separates online learnability from efficient online learnability as it does in the PAC setting (\Cref{ch:pac}).
\end{openprob}

\begin{openprob}[The optimal regret constant, \textsc{ku}]
\label{op:regret-constant}
Ben-David--P\'al--Shalev-Shwartz~\cite{BenDavidPalShalevShwartz2009} give minimax regret $\Theta(\sqrt{\Ldim(\mathcal{H}) \cdot T})$ in the agnostic setting (\cref{rmk:regret-ldim}).  Determine the exact leading constant, and the finite-$T$ correction, matching upper and lower bounds in terms of the Littlestone dimension alone.
\end{openprob}

\begin{openprob}[The multiclass online characterization, \textsc{uk}]
\label{op:online-multiclass}
The mistake tree and $\SOA$ are defined for binary labels.  Determine the multiclass online characterization: which tree dimension (a multiclass Littlestone dimension over $|Y| > 2$) governs the optimal mistake bound, and whether the clean $\mathrm{Opt} = \Ldim$ identity of \cref{thm:littlestone-characterization} survives, paralleling the multiclass PAC question of \Cref{ch:extensions}.
\end{openprob}

\begin{openprob}[Online learnability and privacy, \textsc{ku}]
\label{op:private-online}
Alon--Livni--Malliaris--Moran established that a class is online learnable if and only if it is differentially-privately PAC learnable, with finite Littlestone dimension on both sides.  Formalize this equivalence on top of the verified Littlestone characterization (\leanname{FLT_Proofs.Theorem.Online.optimal_mistake_bound_eq_ldim}), and determine the quantitative dependence of the private sample complexity on $\Ldim(\mathcal{H})$.
\end{openprob}

\begin{openprob}[Quantitative route stability, \textsc{uk}]
\label{op:route-stability-quant}
\Cref{rmk:route-stability} asserts stability away from score ties.  Quantify it: bound the measure of the score-tie boundary set, and the perturbation margin that preserves the route, as functions of the inverse temperature, exhibiting the trade-off between routing sharpness (which improves the Hedge approximation) and stability (which preserves the inherited regret guarantee).
\end{openprob}

\begin{openprob}[Bandit attention, \textsc{ku}]
\label{op:bandit-attention}
In the partial-information (bandit) setting only the loss of the routed key is observed, not the full label.  Determine the regret of online attention routing under bandit feedback (the $\sqrt{kT}$-versus-$\sqrt{T\ln k}$ gap between bandit and full-information exponential weights) and whether route stability mitigates the exploration cost.
\end{openprob}

\begin{openprob}[Halving versus the $\SOA$, \textsc{ku}]
\label{op:halving-soa}
\Cref{rmk:soa-halving} contrasts cardinality with Littlestone dimension.  Characterize the classes on which Halving is sub-optimal by an unbounded factor, where $\log_2|\mathcal{H}|$ (or its effective analogue) exceeds $\Ldim(\mathcal{H})$ arbitrarily, and give the structural property of $\mathcal{H}$ that the $\SOA$ exploits and Halving does not.
\end{openprob}

\begin{openprob}[Online-to-batch for attention, \textsc{uk}]
\label{op:online-to-batch-attention}
The implication online $\Rightarrow$ PAC is verified (\leanname{FLT_Proofs.Theorem.Separation.online_imp_pac}).  Make it quantitative for attention: convert the online mistake bound of a finite-Littlestone attention class into a PAC sample-complexity bound via online-to-batch, and decide whether the conversion is tight or loses a factor, sharpening the relation between \cref{op:ldim-attention} here and the VC question \cref{op:attention-vc} of \Cref{ch:pac}.
\end{openprob}


\subsection*{Counterfactual exercises}

The characterization $\mathrm{Opt}(\mathcal{H})=\Ldim(\mathcal{H})$ rests on several interlocking assumptions. Each exercise below removes or modifies one feature of the game and asks what survives. Each exercise asks you to locate which piece of the argument a given assumption is carrying.

\begin{enumerate}[leftmargin=*,itemsep=3pt]
\item \textbf{The value is two-sided.} The equality $\mathrm{Opt}=\Ldim$ is a conjunction of two theorems: the Standard Optimal Algorithm makes at most $\Ldim$ mistakes (\leanname{soa_mistakes_bounded}) and the adversary forces at least $\Ldim$ (\leanname{adversary_lower_bound}). Show that either half alone gives only a one-sided bound, and that the matching of the two is what makes the value exact.

\item \textbf{Drop realizability.} Remove the constraint that some hypothesis fits every label. The mistake bound becomes vacuous and the right benchmark is regret against the best fixed hypothesis. Show why the agnostic value calls for the regret framework rather than the Littlestone characterization (\cref{op:regret-constant}).

\item \textbf{The gap to PAC.} Online learnability implies PAC learnability but the converse fails. Exhibit the witness: thresholds on $\R$ have $\VC=1$ and $\Ldim=\infty$ (\leanname{online_strictly_stronger_pac}, \leanname{pac_not_implies_online}). Conclude that an adaptive adversary is strictly stronger than any random distribution.

\item \textbf{The conservative leg.} A finite Littlestone bound forces a finite VC bound, so online learnability gives PAC learnability and adds no statistical power (\leanname{online_imp_pac}). Prove this one-way half of the separation directly from $\VC\le\Ldim$.

\item \textbf{An efficient optimal learner.} The $\SOA$ achieves the value information-theoretically, but computing the version-space Littlestone dimensions it needs may be intractable. Investigate which classes admit a polynomial-time-per-round optimal online learner (\cref{op:soa-computable}).

\item \textbf{Depth exceeds set-size.} On thresholds on $\R$ the adversary bisects the surviving interval forever, so the tree dimension is infinite while the shattered-set dimension is one. Show that $\Ldim(\mathcal{H})$ can exceed $\VC(\mathcal{H})$ without bound, and that the slack measures adaptivity (\leanname{ldim_threshold_top}).

\item \textbf{The embedding $\VC\le\Ldim$.} A shattered set of size $d$ becomes a non-adaptive mistake tree of depth $d$ by presenting the same instance at every node of a given depth. Verify that this map preserves size as depth and validity as shattering, so $\VC(\mathcal{H})\le\Ldim(\mathcal{H})$ always (\leanname{BranchWiseLittlestoneDim_ge_VCDim}).

\item \textbf{Multiclass labels.} Replace binary labels by a finite label set, giving each internal node one child per label. Conjecture a multiclass tree-dimension pinning the optimal mistake bound, and determine whether the clean $\mathrm{Opt}=\Ldim$ identity survives (\cref{op:online-multiclass}).
\end{enumerate}
